\section{Formal Model}
\label{sec:model-ra}

The model abstracts away the OS executor, the canonical JSON encoder,
and the receipt ledger; those lower layers are substrate-level
obligations outside the model, discharged by code citation rather than
by mechanized refinement. The institutional core that remains is the
executive-action type, its TTL invariant, the rollback receipt, the
TTL-bounded amendment chain, and the composition theorem connecting
the response side to the parent paper's amendment side.

\paragraph{Executive action.}
An executive action is a structure that cannot be constructed without
a positive TTL witness:
\[
  \mathsf{ExecutiveAction} =
  (\mathit{rid},\,\mathit{t_{\text{issued}}},\,\mathit{ttl},\,
   \mathit{ttlPos},\,\mathit{kind},\,\mathit{rb})
\]
where $\mathit{rid}$ is the receipt identifier, $\mathit{t_{\text{issued}}}$
is the issued-at instant, $\mathit{ttl}$ is the duration in counter
units, $\mathit{ttlPos}$ is the field-level proof $0 < \mathit{ttl}$,
$\mathit{kind}$ is the action class tag, and $\mathit{rb}$ is an
optional rollback receipt. The expiry $\mathit{t_{\text{expires}}} =
\mathit{t_{\text{issued}}} + \mathit{ttl}$ is derived structurally.

The action is closed at instant $t$ iff a rollback receipt exists with
$\mathit{rolledBackAt} \le t$ or the TTL window has elapsed by $t$.
This is recorded as the definitional bridge
\thm{rollback_closes_or_ttl_window_active}, which discharges by a
case split on the TTL window. The bridge is
not the headline theorem; it documents the relationship between the
type-level fields and the propositional closure claim, in the same
role \thm{amendment_admissible_iff_backward_refinement} plays for the
amendment side of the parent paper.

\paragraph{TTL-bounded amendment.}
A TTL-bounded amendment $A$ is a constitutional delta $\delta$ paired
with an issued-at instant $\mathit{t_{\text{issued}}}$ and a positive
TTL duration $\mathit{ttl}$. The active constitution at instant $t$ is
\[
  \mathsf{activeAt}(A, t) =
  \begin{cases}
    \delta.\mathit{old} & \text{if } \mathit{t_{\text{issued}}} + \mathit{ttl} \le t, \\
    \delta.\mathit{new} & \text{otherwise.}
  \end{cases}
\]
The case split is a substrate-level abstraction of the TTL scheduler's
behavior. The scheduler advances the counter and reverts the amendment
at $\mathit{t_{\text{expires}}}$; the model represents this as a
conditional on $t$. Figure~\ref{fig:ttl-window} illustrates the
evolution of an admitted amendment across the window boundary.

\begin{figure}[t]
\centering
\begin{tikzpicture}[
  >={Stealth[length=2.2mm,width=1.6mm]},
  font=\small,
  every node/.style={inner sep=2pt},
]
% main timeline
\draw[->, thick] (0,0) -- (10.5,0) node[right] {$t$};

% tick marks
\draw[thick] (1.2,0.12) -- (1.2,-0.12);
\draw[thick] (7.6,0.12) -- (7.6,-0.12);

% tick labels
\node[below=2pt] at (1.2,-0.12) {$t_0$};
\node[below=2pt] at (7.6,-0.12) {$t_0 + T$};

% region shading
\fill[blue!8] (1.2,0.18) rectangle (7.6,1.05);
\fill[gray!12] (7.6,0.18) rectangle (10.3,1.05);

% region labels
\node[align=center] at (4.4,0.62) {\textbf{in-window state}\\[1pt]
  $\mathsf{activeAt}(A,t)=\delta.\mathit{new}$\\[1pt]
  rollback witness valid};
\node[align=center] at (8.95,0.62) {\textbf{expired}\\[1pt]
  $\mathsf{activeAt}(A,t)=\delta.\mathit{old}$\\[1pt]
  baseline restored};

% admission event arrow
\draw[->, thick] (1.2,1.55) -- (1.2,1.12);
\node[above, align=center] at (1.2,1.55) {action admitted\\[1pt] ${\rm TTL}=T$};

% auto-revert arrow
\draw[->, thick] (7.6,1.55) -- (7.6,1.12);
\node[above, align=center] at (7.6,1.55) {auto-revert\\[1pt] if no rollback receipt};

% brace under in-window
\draw[decorate, decoration={brace, mirror, amplitude=4pt}]
  (1.2,-0.55) -- (7.6,-0.55)
  node[midway, below=6pt, align=center]
  {per-step $\mathsf{BackwardRefines}$ witness preserves baseline on $D$\\[1pt]
   (Theorem~3: baseline preserved on $D$ at every $t$)};
\end{tikzpicture}
\caption{TTL-window evolution. An admitted action is bound to a
positive TTL $T$ at admission time $t_0$. Within the window
$[t_0,\, t_0+T]$, a per-step refinement witness preserves the
baseline. At $t_0+T$, auto-reversion fires if no explicit rollback
receipt has closed the chain. The composition theorem
(Theorem~3) states that the baseline is preserved on the finite
admission domain $D$ at every instant.}
\label{fig:ttl-window}
\end{figure}

\paragraph{The composition theorem.}
The headline composition theorem
\thm{ttl_bounded_amendment_chain_preserves_baseline} states that a
chain of TTL-bounded amendments, each carrying two witnesses over one
explicit finite admission domain $D$ (one that the pre-amendment
constitution backward-refines a fixed baseline on $D$, and one per-step
witness that the post-amendment constitution backward-refines the
pre-amendment constitution on $D$), preserves the baseline on $D$ at
every instant:
\[
  \forall t,\, \forall A \in \mathit{chain},\
    \mathsf{BackwardRefines}\bigl(\mathsf{activeAt}(A, t),\, \mathit{baseline},\, D\bigr).
\]
Here $\mathsf{BackwardRefines}(c', c, D)$ is the substrate's
domain-scoped refinement: every admission view in $D$ that $c'$ admits
is admitted by $c$. The statement is finite-domain, matching the
witness the substrate's amendment check carries; it says nothing about
views outside $D$. The conclusion is pointwise over $A \in \mathit{chain}$
under a universal quantifier, not inductive over an evolving state. For
each amendment $A$, the proof case-splits on the $\mathsf{activeAt}$
conditional. Outside the window,
$\mathsf{activeAt}(A, t) = \delta.\mathit{old}$, and the conclusion is
the post-expiry hypothesis
$\mathsf{BackwardRefines}(\delta.\mathit{old}, \mathit{baseline}, D)$
applied directly. Inside the window,
$\mathsf{activeAt}(A, t) = \delta.\mathit{new}$, and the conclusion
composes the per-step witness
$\mathsf{BackwardRefines}(\delta.\mathit{new}, \delta.\mathit{old}, D)$
with the post-expiry hypothesis: a view in $D$ admitted by
$\delta.\mathit{new}$ is admitted by $\delta.\mathit{old}$ and hence by
the baseline. Both arms consume the per-step witnesses as data rather
than discharging by reflexivity, so each amendment in the chain
contributes substantive content.

The shape is the response-side counterpart of the parent paper's
\thm{essential_preserved_chain}, which composes per-step essential-
admission witnesses over an amendment sequence. The parent theorem
ranges over receipts and a chain of constitutions and is structurally
inductive over the chain; this theorem ranges over instants and a
chain of TTL-bounded amendments and discharges pointwise per amendment
with a case split on the TTL window. The witness-threading pattern is
shared: each step in either chain contributes a per-step refinement
witness that the discharge consumes as data.

\paragraph{Rollback admissibility.}
The supporting reduction \thm{rollback_admission_composes_with_refinement}
states that a rollback receipt admissible under the post-amendment
syntactic constitution remains admissible under the pre-amendment
syntactic constitution given the global syntactic refinement witness
over $\mathsf{admits}$, where $\mathit{rb}$ is the rollback receipt's
admission view:
\[
  \begin{aligned}
    &\mathsf{admits}(c_{\text{new}}, \mathit{rb}) = \text{true}
      \;\wedge\; \bigl(\forall v,\,
         \mathsf{admits}(c_{\text{new}}, v) = \text{true}
         \Rightarrow \mathsf{admits}(c_{\text{old}}, v) = \text{true}\bigr) \\
    &\qquad \Rightarrow\; \mathsf{admits}(c_{\text{old}}, \mathit{rb}) = \text{true}.
  \end{aligned}
\]
The proof is a direct application of the refinement hypothesis to the
rollback receipt's admission. The reduction supplies the rollback-
preservation step the headline composition leans on: a rollback
receipt produced inside an amendment's TTL window remains admittable
when the window expires and the baseline returns.

The two theorems range over the same objects. Both state admission of
an $\mathsf{AdmissionView}$ under a $\mathsf{SyntacticConstitution}$,
decided by $\mathsf{admits}$. What differs is the scope of the
refinement premise. The composition theorem consumes the domain-scoped
$\mathsf{BackwardRefines}(c_{\text{new}}, c_{\text{old}}, D)$, which
quantifies only over views in $D$ and is the witness a checked
$\mathsf{ConstitutionalDelta}$ carries; the rollback reduction consumes
the unrestricted $\mathsf{GlobalSynBackwardRefines}$, which quantifies
over every view. Composing the two therefore carries a side condition:
the rollback receipt's admission view must lie in the delta's domain
$D$, since outside $D$ the composition theorem asserts nothing.

The closure representation the substrate retains under
$\mathsf{Chio.Treaty.Legacy}$ is related to the syntactic one
pointwise: the closure induced by a syntactic constitution admits a
view exactly when $\mathsf{admits}$ does. That is the parent
substrate's $\mathsf{bridge\_pointwise}$, restated in the paper-local
file as \thm{closure_to_syntactic_admission_bridge} so that both
reductions can be read against either representation. It is an
instance of the parent theorem, not a new proof.

\paragraph{Bilateral destructive admission.}
The bilateral reduction \thm{destructive_action_requires_bilateral_admission}
specializes the parent paper's
\thm{treaty_admission_iff_predicate_intersection} to the destructive
subclass of the action taxonomy. For a bilateral envelope binding one
admission view to a destructive action class under a bilateral treaty,
admission holds iff both the device polity's and the operator polity's
predicates accept that view. The reduction is structural rather
than novel: the parent theorem already establishes the bilateral
intersection; this theorem renames it for the destructive class. The
contribution at this point is the typed envelope, not a new theorem.

\paragraph{Threat model.}
The adversary can author and sign requests, attempt to skip the
rollback receipt, forge or delete the rollback receipt in the local
ledger, backdate or freeze the substrate counter, omit the
acknowledgement receipt, present a rollback receipt for an execution
it did not perform, or claim that an inverse executor completed when
it did not. The adversary cannot break Ed25519, SHA-256 collision
resistance, or canonical JSON correctness. The substrate counter
itself is an assumption: a compromised counter falsifies the TTL claim
in the same way a compromised signing key falsifies admission. The
counter is named as a foundational assumption in
Section~\ref{sec:limitations-ra}.

The model discharges two obligations. First, the typed executive-
action constructor refuses to build a witness without a positive TTL
and a class tag, so the type system rules out unbounded acts on the
construction side. Second, the composition theorem pins the chain's
trajectory to the baseline at every instant, so an amendment ratchet
that aggregates predicate drift through a sequence of TTL-bounded
amendments cannot accumulate without violating a per-step refinement
witness. The third obligation, namely that the rollback receipt
faithfully witnesses the inverse executor's completion, is a runtime
obligation discharged by the executor's signed completion against the
forward effect's content hash and is not a theorem in the bounded
model.
